\section{Small central splittings from finite-dimensional models}
\label{n:sec:small}

We now construct a positive contraction that approximately splits each
coefficient trace in a small proportion $\lambda$.  The trace and
projection errors will be of order $\lambda^{3/2}$, small enough to
vanish under the amplification in Section~\ref{n:sec:amplification}.
We first state the output, then construct finite-dimensional models and
control the sum of projections selected to commute with them.

\begin{proposition}\label{n:prop:small}
Let $0<\lambda<1$, let $\mathcal F$ be a finite set of self-adjoint
contractions in $A$, and let $\rho>0$. There is a positive contraction
$e\in A$ such that, for $\sigma\in T$ and $x\in\mathcal F$,
\begin{equation}\label{n:eq:small}
 \begin{split}
 \|[e,x]\|_{2,\sigma}&<\rho,\\
 |\sigma(ex)-\lambda\sigma(x)|&<\lambda^{3/2}+\rho,\\
 \sigma(e-e^2)&<\lambda^{3/2}+\rho.
 \end{split}
\end{equation}
\end{proposition}

\subsection{Averaged finite-dimensional models}\label{n:sec:models}

We now approximate a finite tuple from $A$ by finite-dimensional
subalgebras of $D$. Hyperfiniteness first supplies an order-zero model at
each trace. Convex averaging makes their weighted squared error uniformly
small, and projection selection replaces each order-zero model by a
finite-dimensional subalgebra of $D$. The number of models may grow as
the accuracy improves; only their weighted total error will enter the
next construction.

\begin{proposition}\label{n:prop:models}
For a finite set $\mathcal F$ of self-adjoint contractions in $A$ and
$\Delta>0$, there are unital finite-dimensional subalgebras $D_i\subset D$,
self-adjoint contractions $x^{(i)}\in D_i$ for each $x\in\mathcal F$, and
positive weights $w_i$ summing to one, such that
\begin{equation}\label{n:eq:models}
 E_i=\sum_{x\in\mathcal F}(x-x^{(i)})^2,
 \qquad \sum_iw_i\tau(E_i)<\Delta\quad(\tau\in \Lambda).
\end{equation}
\end{proposition}
\begin{proof}
Consider all finite-dimensional unital $E^0$, order-zero maps
$\theta:E^0\to A$, and self-adjoint contractions $b_x\in E^0$. Their
error functions
\[
 L(\sigma)=1-\sigma(\theta(1))+
       \sum_{x\in\mathcal F}\sigma((x-\theta(b_x))^2)
\]
are nonnegative continuous affine functions on $T$. At each fixed
$\sigma$, hyperfiniteness of $M_\sigma$ supplies a unital
finite-dimensional subalgebra approximating the tuple in $2$-norm by
self-adjoint contractions. Lift its inclusion through
\eqref{n:eq:kaplansky} and choose a coordinate. This makes $L(\sigma)$
arbitrarily small, including its unit-defect term.

For every $\kappa>0$, a finite convex combination of these functions has
maximum less than $\kappa$. Otherwise separate their convex hull from the
open convex set $\{f\in C(T,\R):\max_Tf<\kappa\}$. The separating
functional is positive, since this set is closed under adding arbitrary
negative functions. Normalize it to a probability measure $\mu$.
Its barycenter $\sigma_\mu$ belongs to compact convex $T$, and affinity
gives $L(\sigma_\mu)=\int L\,d\mu\ge\kappa$ for every candidate $L$,
a contradiction. Denote the resulting data by $D_i^0,\theta_i,b_{i,x},w_i$.

We next replace each $\theta=\theta_i$, now viewed in $D$, by a
homomorphism on a projection, controlling the loss of unit trace by
the original unit defect. Write $\theta=h\pi$, and put
$f(s)=(2s-1)_+$, $g(s)=\min(1,2s)$. For the $e_{11}$ corner of a matrix
summand, apply Proposition~\ref{n:prop:selection} to
\[
 (g(\theta(e_{11})),f(\theta(e_{11})),f(\theta(e_{11}))).
\]
The support condition holds because $gf=f$, and the trace inequality
is equality. We obtain a projection $v$ with
\[
 g(\theta(e_{11}))v=v,\qquad
 \tau(v)=\tau(f(\theta(e_{11})))\quad(\tau\in \Lambda).
\]
It lies under $\pi(e_{11})$. The elements $\pi(e_{j1})v$ belong to $D$,
because
\[
 \pi(e_{j1})g(\theta(e_{11}))=\pi(e_{j1})g(h)\in D.
\]
Repeating $v$ along each matrix diagonal and summing over the summands
gives a projection $Q\in D$ commuting with $\pi(D_i^0)$, with
$Q\pi(d)\in D$ for $d\in D_i^0$. Traciality and repetition give
\[
 \tau(Q)=\tau(f(h))\quad(\tau\in \Lambda).
\]
Set $D_i=\C(1-Q)+Q\pi(D_i^0)$ and $x^{(i)}=Q\pi(b_{i,x})$.
This is a unital finite-dimensional subalgebra; zero image summands are
simply absent. No commutation of $Q$ with $h$ is needed. For a contraction
$b\in D_i^0$, tracial $L^2$ contractivity and the triangle inequality give
\begin{align*}
 \|\theta(b)-Q\pi(b)\|_{2,\tau}
 &\le\|1-h\|_{2,\tau}+\|1-Q\|_{2,\tau}\\
 &\le(1+\sqrt2)\tau(1-h)^{1/2},
\end{align*}
since $1-f(s)\le2(1-s)$. With $m=|\mathcal F|$ and
$C_0=(1+\sqrt2)^2$, this implies
\[
 \tau(E_i)\le2\sum_{x\in\mathcal F}
       \tau((x-\theta_i(b_{i,x}))^2)+2mC_0\tau(1-\theta_i(1)).
\]
The weighted sum is at most $(2+2mC_0)\kappa$. Take $\kappa$ sufficiently
small. Limit traces restrict to $T$ by Lemma~\ref{n:lem:unit}, so the bound
holds on $\Lambda$.
\end{proof}

\subsection{Truncation and the combined cost}

We will use Corollary~\ref{n:cor:cost} to select one projection for each
model, commuting with that model. Its cost will control both its overlap
with earlier choices and its error against the prescribed coefficients.
The sum need not be a contraction. The following estimate permits truncation
while retaining bounds independent of the number of models.

\begin{lemma}\label{n:lem:truncation}
Let $p_1,\ldots,p_N$ be projections in a unital $C^*$-algebra with a
tracial state $\tau$. Put
\[
 s=\sum_ip_i,\qquad e=\min(s,1),\qquad
 \lambda=\sum_i\tau(p_i),\qquad R=\sum_i\tau\Bigl(p_i\sum_{j<i}p_j\Bigr).
\]
For every self-adjoint $x$,
\begin{equation}\label{n:eq:trunc-comm}
 \|[e,x]\|_{2,\tau}
 \le2\Bigl(2\sum_i\|[p_i,x]\|_{2,\tau}^2\Bigr)^{1/2}.
\end{equation}
Moreover,
\begin{equation}\label{n:eq:trunc-trace}
 \tau(s-e),\ \tau(e-e^2)\le\sqrt{2\lambda R}.
\end{equation}
\end{lemma}
\begin{proof}
Use the unnormalized matrix trace and set
\[
 W=(p_1,\ldots,p_N),\quad
 V=\begin{pmatrix}0&W\\W^*&0\end{pmatrix},\quad
 Z=\operatorname{diag}(x,\ldots,x).
\]
For $l(r)=\min(r^2,1)$, the first diagonal corner of $l(V)$ is $e$.
The function $l$ is globally 2-Lipschitz. Left and right multiplication
by $V$ are commuting self-adjoint operators on tracial $L^2$, so their
joint spectral calculus gives
\[
 \|[l(V),Z]\|_2\le2\|[V,Z]\|_2.
\]
The squared norm on the right before the factor $4$ is
$2\sum_i\|[p_i,x]\|_{2,\tau}^2$. Compressing the left side to the first
corner proves \eqref{n:eq:trunc-comm}. This argument uses scalar
Lipschitzness on $L^2$, not operator-norm Lipschitzness.

Let $S$ be the support of $s$ in the tracial von Neumann representation.
The row $W$ is supported on the right by
$\operatorname{diag}(p_1,\ldots,p_N)$, so its polar decomposition gives
$\tau(S)\le\lambda$. Since $\tau(s^2)=\lambda+2R$,
\[
 \tau((s-S)^2)=2R-\lambda+\tau(S)\le2R.
\]
Both $s-e$ and $e-e^2$ are positive, supported under $S$, and bounded
above by $|s-S|$. Cauchy--Schwarz gives
\[
 \tau(|s-S|)\le\tau(S)^{1/2}\tau((s-S)^2)^{1/2}
       \le\sqrt{2\lambda R},
\]
which proves \eqref{n:eq:trunc-trace}.
\end{proof}

\begin{proof}[Proof of Proposition~\ref{n:prop:small}]
We first work in $D$. Choose a large parameter $K>1$, and apply
Proposition~\ref{n:prop:models} with $\Delta=K^{-2}$. Recursively use
Corollary~\ref{n:cor:cost} to choose $p_i\in D\cap D_i'$ with
\[
 t_i=\lambda w_i,\qquad Y_i=s_{i-1}+KE_i,\qquad
 s_j=\sum_{i\le j}p_i,
\]
and tolerance $K^{-1}$. In particular,
$\tau(p_i)=t_i$ and $\tau(p_ix^{(i)})=t_i\tau(x^{(i)})$.
For each $\tau\in \Lambda$, write
\[
 R_\tau=\sum_i\tau(p_is_{i-1}),\qquad
 H_\tau=\sum_i\tau(p_iE_i).
\]
Both are nonnegative, and the single selected cost controls both:
\begin{align}
 R_\tau+KH_\tau
 &\le\sum_it_i\Bigl(\sum_{j<i}t_j+K\tau(E_i)+K^{-1}\Bigr)
 \le\lambda^2/2+2\lambda/K.\label{n:eq:joint-cost}
\end{align}
In particular $H_\tau\le\lambda^2/(2K)+2\lambda/K^2$, uniformly in
$\tau$ and independently of the number of models.

Since $[p_i,x^{(i)}]=0$, for $x\in\mathcal F$,
\[
 \sum_i\|[p_i,x]\|_{2,\tau}^2\le4H_\tau.
\]
The positive functional $z\mapsto\tau(p_i z)$ has norm $t_i$, whence
\[
 |\tau(p_i(x-x^{(i)}))|\le\sqrt{t_i\tau(p_iE_i)}.
\]
Scalar Cauchy--Schwarz and \eqref{n:eq:models} therefore give, for $s=s_N$,
\begin{equation}\label{n:eq:sum-moment}
 |\tau(sx)-\lambda\tau(x)|\le\sqrt{\lambda H_\tau}+\lambda/K.
\end{equation}

Set $e_D=\min(s,1)$. Lemma~\ref{n:lem:truncation} gives
\[
 \|[e_D,x]\|_{2,\tau}\le2\sqrt{8H_\tau},\qquad
 \tau(s-e_D),\ \tau(e_D-e_D^2)\le\sqrt{2\lambda R_\tau}.
\]
Replacing $s$ by $e_D$ in \eqref{n:eq:sum-moment} costs at most
$\tau(s-e_D)$. By \eqref{n:eq:joint-cost}, the upper bounds for the trace and projection errors
tend to $\lambda^{3/2}$ as $K\to\infty$, while the commutator bound
tends to zero. Choose $K$ giving the desired inequalities with slack
strictly smaller than $\rho$. A positive contraction lift of $e_D$ and
one good coordinate then give \eqref{n:eq:small} in $A$.
\end{proof}
